\chapter{Fundamental types, constants and variables}
\label{chap:types}

Ymir is an imperative language, and an imperative program reads like a recipe:
one instruction after another, in the order written, each leaving the program a
little further along than it found it. What the instructions work on is data, and
data is held in constants and variables, each of a particular type.

Statements come in two kinds. A declaration brings a new name into existence and
says what it holds. An instruction changes something -- the contents of a
variable, or the state of the program at large. This chapter covers both, and the
types that the values behind those names can have.

\minitoc%

\section{Constant and variables}
\label{sec:types:constants_variables}

Every symbol a program declares has a name, or \textit{identifier}. Identifiers
are case-sensitive -- \token{count} and \token{Count} are two different names --
and they are built from letters, digits and underscores, subject to two rules: an
identifier contains at least one letter, and no digit may appear before the first
letter. Any number of leading underscores is allowed.

So \token{__id1__} is a valid identifier, while \token{__1id__} is not: the digit
arrives before any letter does.

\notebox{
  The formal grammar of an identifier is the following: \token{(\_)* [a-zA-Z]([a-zA-Z]|[0-9]|\_)*}
}

A handful of names obey those rules and are still unavailable: the language has
reserved them for statements and operations of its own. They are the
\textit{keywords}, and there are no others than these.

\token{alias}, \token{assert}, \token{async}, \token{atomic}, \token{await},
\token{break}, \token{cast}, \token{catch}, \token{class}, \token{const},
\token{continue}, \token{copy}, \token{cte}, \token{dcopy}, \token{def},
\token{dg}, \token{dmut}, \token{else}, \token{entity}, \token{enum},
\token{expand}, \token{extern}, \token{false}, \token{fn}, \token{for},
\token{future}, \token{if}, \token{impl}, \token{in}, \token{is}, \token{lazy},
\token{let}, \token{loop}, \token{macro}, \token{match}, \token{mod},
\token{move}, \token{mut}, \token{none}, \token{null}, \token{of},
\token{panic}, \token{\_\_pragma}, \token{record}, \token{ref}, \token{return},
\token{self}, \token{spawn}, \token{static}, \token{super}, \token{template},
\token{\_\_test}, \token{throw}, \token{throws}, \token{trait}, \token{true},
\token{typeof}, \token{union}, \token{unsafe}, \token{use}, \token{\_\_version},
\token{while}, \token{with}, \token{yield}, and the wildcard \token{\_}.


Each of them is introduced where the book needs it; none is used before it has
been explained.

\vfill\pagebreak
\section{Variable declaration}
\label{sec:types:variable_declaration}

Inside a function, \token{let} declares a variable: an identifier that stands for
a value. Every value has a type, and the type is what tells the program how to
read the bits and what may be done with them.

A useful picture is a \textit{box}. The variable is the label on the box, the
value is what is inside, and the type is the shape of the box -- which is what
keeps a value of the wrong form from being put in it. This chapter stays with the
types whose values fit in a single box, the scalar types:

\vspace{-10pt}%
\begin{itemize}
\setlength\itemsep{0.2mm}
\item integer types,
\item floating-point types,
\item character types,
\item the boolean type.
\end{itemize}\vspace{-10pt}%

Each of the first three comes in several variants, covering different ranges of
values. The sections that follow take them one at a time; for now, one of each is
enough -- \token{i32} for integers, \token{f32} for floating-point numbers,
\token{c8} for characters -- alongside \token{bool}.

The program below declares one variable of each: \token{x}, \token{y}, \token{z}
and \token{w}.

\begin{lstlisting}[style=coloredverbatim, caption=simple variable declaration, label=lst:types:basic_variables]
use std::io;

fn main() {
  let x: i32 = 10;
  let y: f32 = 3.1415;
  let z: c8 = 'c';
  let w: bool = false;

  println("X : ", x);
  println("Y : ", y);
  println("Z : ", z);
  println("W : ", w);
}
\end{lstlisting}

\begin{figure}[h]
  \vspace{-5pt}%
  \centering
  \adjustbox{scale=0.90, max width=\linewidth}{
    \begin{tikzpicture}

      \draw (0, 0) -- coordinate (X) (0, 0);
      \draw (4, 0) -- coordinate (Y) (4, 0);
      \draw (8, 0) -- coordinate (Z) (8, 0);
      \draw (12, 0) -- coordinate (W) (12, 0);

      \fill[olive!10] ($(X)$) rectangle ($(X)+(3, 5)$);
      \filldraw ($(X)+(1.5, -0.2)$) node [align=center, below] {\textit{x}};
      \filldraw ($(X)+(1.5, 2.5)$) node [align=center] {10};
      \filldraw ($(X)+(-0.4, 2.5)$) node [align=center, rotate=90] {\small {\textit{32 bits}}};
      \draw[Stealth-Stealth] ($(X)+(-0.1, 0.1)$) -- ($(X)+(-0.1, 4.9)$);

      \fill[green!10] ($(Y)$) rectangle ($(Y)+(3, 5)$);
      \filldraw ($(Y)+(1.5, -0.2)$) node [align=center, below] {\textit{y}};
      \filldraw ($(Y)+(-0.4, 2.5)$) node [align=center, rotate=90] {\small {\textit{32 bits}}};
      \filldraw ($(Y)+(1.5, 2.5)$) node [align=center] {3.1415};
      \draw[Stealth-Stealth] ($(Y)+(-0.1, 0.1)$) -- ($(Y)+(-0.1, 4.9)$);

      \fill[cyan!10] ($(Z)$) rectangle ($(Z)+(3, 1.25)$);
      \filldraw ($(Z)+(1.5, -0.2)$) node [align=center, below] {\textit{z}};
      \filldraw ($(Z)+(-0.4, 0.625)$) node [align=center, rotate=90] {\small {\textit{8 bits}}};
      \filldraw ($(Z)+(1.5, 0.625)$) node [align=center] {'c'};
      \draw[Stealth-Stealth] ($(Z)+(-0.1, 0.1)$) -- ($(Z)+(-0.1, 1.15)$);

      \fill[purple!10] ($(W)$) rectangle ($(W)+(3, 1.25)$);
      \filldraw ($(W)+(1.5, -0.2)$) node [align=center, below] {\textit{w}};
      \filldraw ($(W)+(1.5, 0.625)$) node [align=center] {false};
      \filldraw ($(W)+(-0.4, 0.625)$) node [align=center, rotate=90] {\small {\textit{8 bits}}};
      \draw[Stealth-Stealth] ($(W)+(-0.1, 0.1)$) -- ($(W)+(-0.1, 1.15)$);

    \end{tikzpicture}
  }
  \caption{\label{fig:types:variable_bins} Representation of variables of Listing~\ref{lst:types:basic_variables} as boxes}
\end{figure}


Entering \token{main} creates the four boxes of
Figure~\ref{fig:types:variable_bins}, each holding its value. Naming a box hands
back what is inside it, which is what the four \token{println} calls on lines~9
to~12 do. The boxes are a picture and nothing more -- no such object exists at
run time -- but it is the right picture for everything in this chapter.

A box that is never opened is a mistake, and the compiler treats it as one: a
variable that is declared and never read stops the compilation, with the message
that the symbol ``was declared but never used''. That is why the listings of this
chapter end by printing what they declare.

The type sits between \token{:} and \token{=} in the declaration, and it is
optional: the compiler infers the shape of the box from the first value put in
it. What is not optional is that the shape, once fixed, stays fixed; only values
of that shape go in afterwards.

The four boxes above are also sealed. Their contents cannot be replaced, which
makes them constants. \token{mut}, written after \token{let}, produces a box that
can be opened again -- a \textit{mutable} variable. The book calls both kinds
variables, as everyone does, though strictly it is only the mutable ones that
vary.



Replacing the contents of a mutable variable is the job of \token{=}. The
assignment holds only if both sides have the same type.

The two sides have names. The left operand is the \textit{lvalue}, the right one
the \textit{rvalue} -- \textit{left value} and \textit{right value}. In the box
picture the lvalue is the box and the rvalue is what goes into it. In
Listing~\ref{lst:types:value_change} the lvalue is \token{x}, twice over; the
output follows the listing.

\begin{lstlisting}[style=coloredverbatim, caption=Example of mutable variable, label=lst:types:value_change]
use std::io;

fn main() {
  let mut x = 10;

  x = 42;
  println("X : ", x);

  x = 314;
  println("X : ", x);
}
\end{lstlisting}
\vspace{-10pt}%
\begin{lstlisting}[style=bashVerb, label=lst:types:value_change_output]
X : 42
X : 314
\end{lstlisting}

A second picture explains the run: a cursor travelling down the instructions,
starting at \token{main}.

It creates a box for \token{x} holding \token{10}. One line down, the assignment
puts \token{42} in that box instead. One line down again, \token{println} reads
the box and writes what it finds. Lines~9 and~10 repeat the pair.

That is all imperative programming is -- state that instructions change, one
after another, as control moves through them -- and procedural programming is the
name for doing it with functions. Ymir offers several ways to organise a program,
and all of them rest on this.

\vfill\pagebreak%
\section{Operators}
\label{sec:types:operators}

A program that could only move values from box to box would not be worth much.
Operators are what build new values out of old ones, and an rvalue is usually not
a bare value but an expression: one or more variables combined by operators. They
fall into four groups.

\vspace{-5pt}%
\begin{itemize}
  \setlength\itemsep{-4pt}
\item \textbf{Arithmetic} operators -- addition, subtraction, multiplication and
  the rest -- apply to integers and floating-point numbers, and return a value of
  the same type as their operands.
\smallskip{}

\item \textbf{Comparison} operators relate two values and return a
  \token{bool}. They are what a program branches on.
\smallskip{}

\item \textbf{Bitwise} operators work on the individual bits of a value. They
  wait for a later chapter, because using them correctly means knowing first how
  a value is laid out in memory.
\smallskip{}

\item \textbf{Boolean} operators combine \token{bool} values according to
  Boolean logic.
  \smallskip{}

\end{itemize}

The listing below puts the first two groups to work on \token{i32} values.

\begin{lstlisting}[style=coloredverbatim, label=lst:types:operator_example, caption=Example of basic operators on \token{i32} values]
use std::io;

fn main() {
  let x = 12;
  let y = 24;

  let z = x + y;
  let w = z * 4;

  println("Z : ", z);
  println("W : ", w);
  let b = z > x;

  println("Z > X : ", b);
}
\end{lstlisting}
\vspace{-10pt}%
\begin{lstlisting}[style=bashVerb, label=lst:types:operator_example_output]
Z : 36
W : 144
Z > X : true
\end{lstlisting}

Several operators can appear in one expression, and then the question is which
runs first. As in mathematics, each operator has a precedence, and the higher
precedence goes first; Table~\ref{tab:types:operator_precedence} gives the order
for all of them.

Precedence is easiest to see as a tree. Each node is an operator, its children
are the operands it works on, and the leaves are the values written in the source
code. Evaluation runs from the leaves upwards, each node collapsing to a single
value, until the root is all that is left.

Figure~\ref{fig:types:operator_precedence} is the tree for
\token{1 + 3 * (4 - 9)}. Parentheses outrank every operator, so \token{4 - 9}
collapses first, even though \token{-} would otherwise yield to \token{*}.

\begin{table}[h]
  \centering
  \adjustbox{scale=0.9, max width=\linewidth}{
    \begin{tabular}{p{.15\textwidth}|p{.55\textwidth}|p{0.30\textwidth}}
      \toprule[0.6pt]
      \textbf{Precedence} & \textbf{Operators} & \textbf{Usage} \\
      \toprule[0.6pt]
      1 & \token{=} \token{/=} \token{-=} \token{+=} \token{*=} \token{\%=} \colorbox{background!90!gray!60}{\mtilde=} \token{<<=} \token{>>=}& Assignment \\
      2 & \token{||}& Boolean \\
      3 & \token{&&}& Boolean \\
      4 & \token{<} \token{>} \token{<=} \token{>=} \token{\!=} \token{==} \token{of} \token{is} \token{in} \token{\!of} \token{\!is} \token{\!in} \token{<=>} \token{<\!=>} & Comparison \\
      5 & \token {...} \token {..} & Range \\
      6 & \token {<<} \token {>>} & Bitwise shift \\
      7 & \token {|} \token {^} \token {&} & Bitwise \\
      8 & \token {+} \colorbox{background!90!gray!60}{\mtilde} \token {-} & Arithmetics \\
      9 & \token {*} \token {\%} \token {/} & Arithmetics  \\
      10 & \token {\^\^} & Exponent \\
      \midrule[0.2pt]
    \end{tabular}
  }
  \caption{\label{tab:types:operator_precedence} Operator precedence}
\end{table}


Figure~\ref{fig:types:operator_precedence_no_par} is the tree for
\token{1 + 3 * 4 - 9}, the same expression with the parentheses taken out. The
shape is different, and so is the answer: with nothing to override it, precedence
alone decides.

\begin{figure}[h]
  \centering
  \begin{subfigure}[t]{0.46\textwidth}
    \centering
    \begin{tikzpicture}
      \draw (0, 0.2) -- coordinate (A) (0, 0.2);
      \draw (0, -0.8) -- coordinate (B) (0, -0.8);
      \draw (1, -1.6) -- coordinate (C) (1, -1.6);
      \draw (2, -2.4) -- coordinate (D) (2, -2.4);


      \filldraw ($(A)+(0, -0.05)$) node[align=center] {+};
      \filldraw ($(B)+(1,-0.1)$) node[align=center] {*};
      \filldraw ($(B)+(-1,0)$) node[align=center] {1};

      \filldraw ($(C)+(-1,0)$) node[align=center] {3};
      \filldraw ($(C)+(1,-0.1)$) node[align=center] {-};

      \filldraw ($(D)+(-1,0)$) node[align=center] {4};
      \filldraw ($(D)+(1,0)$) node[align=center] {9};

      \draw[-] ($(A)+(0.1, -0.2)$) -- ($(B)+(0.9, 0.2)$);
      \draw[-] ($(A)+(-0.1, -0.2)$) -- ($(B)+(-0.9, 0.2)$);

      \draw[-] ($(B)+(1.1, -0.2)$) -- ($(C)+(0.9, 0.2)$);
      \draw[-] ($(B)+(0.9, -0.2)$) -- ($(C)+(-0.9, 0.2)$);

      \draw[-] ($(C)+(1.1, -0.2)$) -- ($(D)+(0.9, 0.2)$);
      \draw[-] ($(C)+(0.9, -0.2)$) -- ($(D)+(-0.9, 0.2)$);

    \end{tikzpicture}

    \caption{\label{fig:types:operator_precedence} operator precedence for \token{1 + 3 * (4 - 9)}}
  \end{subfigure}%
  \begin{subfigure}[t]{0.46\textwidth}
    \centering
    \begin{tikzpicture}

      \draw (0, 0.2) -- coordinate (A) (0, 0.2);
      \draw (0, -0.8) -- coordinate (B) (0, -0.8);
      \draw (1, -1.6) -- coordinate (C) (1, -1.6);
      \draw (0, -2.4) -- coordinate (D) (0, -2.4);

      \filldraw ($(A)+(0,-0.1)$) node[align=center] {-};
      \filldraw ($(B)+(1,0)$) node[align=center] {9};
      \filldraw ($(B)+(-1,-0.05)$) node[align=center] {+};

      \filldraw ($(C)+(-1,-0.1)$) node[align=center] {*};
      \filldraw ($(C)+(-3,0)$) node[align=center] {1};

      \filldraw ($(D)+(-1,0)$) node[align=center] {3};
      \filldraw ($(D)+(1,0)$) node[align=center] {4};

      \draw[-] ($(A)+(0.1, -0.2)$) -- ($(B)+(0.9, 0.2)$);
      \draw[-] ($(A)+(-0.1, -0.2)$) -- ($(B)+(-0.9, 0.2)$);

      \draw[-] ($(B)+(-0.9, -0.2)$) -- ($(C)+(-1, 0.2)$);
      \draw[-] ($(B)+(-1.1, -0.2)$) -- ($(C)+(-3, 0.2)$);

      \draw[-] ($(C)+(-0.9, -0.2)$) -- ($(D)+(0.9, 0.2)$);
      \draw[-] ($(C)+(-1.1, -0.2)$) -- ($(D)+(-0.9, 0.2)$);

    \end{tikzpicture}

    \caption{\label{fig:types:operator_precedence_no_par} operator precedence for \token{1 + 3 * 4 - 9}}
  \end{subfigure}
  \caption{Example of operator precedence}
\end{figure}


The table also lists the compound assignments, \token{+=}, \token{-=} and their
kin. Each folds an operation into an assignment: the lvalue serves as the left
operand, the rvalue as the right, and the result goes back into the lvalue. So
\token{+=} adds to a variable and stores the sum in one step.

The rvalue is taken whole, whatever its precedence. In \token{x *= y + 1} the
addition binds more loosely than the multiplication, and it still happens first:
the statement means \token{x = x * (y + 1)}, not \token{x = x * y + 1}.

The compiler always knows which operator goes first. The reader of the code may
not, so in a long expression it is worth adding parentheses that change nothing,
or splitting the work over several lines.

\begin{lstlisting}[style=coloredverbatim]
let mut x = 10;
let y = 3;

x *= 2 + y; // rewritten into x = x * (2 + y);
println("X : ", x);
\end{lstlisting}
\vspace{-10pt}%
\begin{lstlisting}[style=bashVerb]
X : 50
\end{lstlisting}


Everything so far has been binary -- two operands. Unary operators take one, and
bind more tightly than any binary operator. Ymir has four: \token{-}, \token{&},
\token{*} and \token{!}. What each one means depends on the type it is applied
to:
\vspace{-8pt}%
\begin{itemize}
\setlength\itemsep{0.2mm}
\item \token{-} on an integer gives its opposite.
\item \token{-} does not apply to a \token{bool} at all.
\item \token{!} on a \token{bool} gives the other \token{bool} -- logical
  \textit{not}.
\item \token{!} on an integer is a bitwise operation, and waits for the chapter
  that covers those.
\end{itemize}


\vfill\pagebreak
\section{Integer types}
\label{sec:types:integers}

Ymir has ten integer types, and their names say what they are. The letter is the
sign: \token{i8}, \token{i16}, \token{i32}, \token{i64} and \token{isize} are
signed and hold negative values as well as positive ones; \token{u8},
\token{u16}, \token{u32}, \token{u64} and \token{usize} are unsigned and hold
only non-negative ones. The number is the width in bits, and the width fixes how
many distinct values the type can tell apart. \token{isize} and \token{usize}
are the exception: their width is whatever the machine works in naturally, 64
bits on ARM64, 32 on x86.

\subsection{Literal representation}
\label{sec:types:integer_literals}

An integer literal is a run of digits written directly in the source code, in
decimal and without spaces -- \token{12} and \token{36} in the listings so far.
Long ones are hard to read, so underscores may be scattered through them as
separators: \token{1\_000\_000\_000\_000} is one trillion, and the underscores
are ignored.

A literal with nothing else attached is an \token{i32}. A suffix names another
type: \token{12u8} is an 8-bit unsigned value, \token{12i64} a 64-bit signed
one.

\notebox{
  More advanced integer literals can be written in octal, hexadecimal,
  or binary form. These representations are covered in detail in the advanced
  chapter on integers (see Section~\ref{sec:scalars:integers}).
}

\subsection{Memory representation}
\label{sec:types:integer_memory}

A computer stores numbers in base~$2$. Each digit is a bit, worth $0$ or $1$, and
each position is worth twice the one to its right: $2^0$ at the far right, then
$2^1$, $2^2$, and so on. An 8-bit integer is thus an eight-digit number in
base~$2$, and occupies one byte.

\begin{table}[h!]
  \centering
  \label{tab:types:integer_ranges}
  %% Columns sized to their contents rather than to fixed fractions of the text
  %% width. As p{} columns, the Type cells were too narrow for their two chips
  %% and the Signed Range cells too narrow for their formula, so both wrapped --
  %% which put each type one line below the row it labels, and broke $2^{N-1}-1$
  %% across two lines after the minus sign.
  \begin{tabular}{l|l|c|c}
    \toprule[0.6pt]
    \textbf{Type} & \textbf{Bits} & \textbf{Signed Range} & \textbf{Unsigned Range} \\
    \toprule[0.6pt]
    \token{i8} / \token{u8}      & 8    & $-2^7 \ldots 2^7-1$        & $0 \ldots 2^8-1$ \\
    \token{i16} / \token{u16}    & 16   & $-2^{15} \ldots 2^{15}-1$  & $0 \ldots 2^{16}-1$ \\
    \token{i32} / \token{u32}    & 32   & $-2^{31} \ldots 2^{31}-1$  & $0 \ldots 2^{32}-1$ \\
    \token{i64} / \token{u64}    & 64   & $-2^{63} \ldots 2^{63}-1$  & $0 \ldots 2^{64}-1$ \\
    \token{isize} / \token{usize} & system-dependent (N) & $-2^{N-1} \ldots 2^{N-1}-1$ & $0 \ldots 2^N-1$ \\
    \midrule[0.2pt]
  \end{tabular}
  \caption{Integer Types and Ranges in Ymir (expressed as powers of 2)}
\end{table}



Unsigned 8-bit integers \token{u8} use all 8 bits to represent non-negative
numbers, ranging from $0$ to $2^8-1 = 255$. Each bit contributes to the total
value depending on its position: the rightmost bit represents 1, the next bit 2,
then 4, 8, 16, 32, 64, and finally 128 for the leftmost bit. For instance, the
number 13 in binary is $00001101$, which corresponds to:
\begin{equation}
  (0\times{}128) +
  (0\times{}64) + (0\times{}32) + (0\times{}16) + (1\times{}8) + (1\times{}4) + (0\times{}2) +
  (1\times{}1) = 13
\end{equation}

A signed \token{i8} has to fit negative numbers into the same eight bits, and
does it by two's complement. The leftmost bit -- the most significant, or MSB --
carries the sign: $0$ for positive, $1$ for negative. Positive values are
unchanged, so $13$ is still $00001101$. A negative value is built in three steps:
write the magnitude in binary ($00001101$), invert every bit ($11110010$), add
one ($11110011$). That last pattern is $-13$.

Spending a bit on the sign costs range but not capacity. An \token{i8} still
distinguishes 256 values, exactly as a \token{u8} does; they are simply a
different 256, from $-2^7$ to $2^7-1$, which is $-128$ to $127$.

\notebox{
  The unary operator \token{!}, when applied to an integer, returns its one’s
  complement—that is, it inverts all the bits of the number. For example, the
  operation \token{(!13) + 1} produces the two’s complement of $13$, which is equal
  to \token{-13}.
}

Nothing above depends on the width being eight. For $N$ bits, an unsigned type
covers $0$ to $2^N-1$ and a signed one $-2^{N-1}$ to $2^{N-1}-1$. Base~$2$ can
write any whole number given enough digits; a machine has a fixed number of them,
which is why every integer type has ends.


\subsection{Integer overflow}
\label{sec:types:integer_overflow}

Those ends are hard. When a calculation lands outside them, the result is not the
number that was wanted -- it is \textit{integer overflow}, and the value wraps
around to the other end of the range.

A \token{u8} runs from $0$ to $255$, so $255 + 1$ is $0$ and $0 - 1$ is $255$. An
\token{i8} runs from $-128$ to $127$, so $127 + 1$ is $-128$. Ymir does not
detect this and does not prevent it. It cannot be caught at compile time, since
the operands are generally not known until the program runs, and checking every
arithmetic operation at run time would cost more than it is worth.

The consequence is worth stating plainly: an overflowed calculation produces a
wrong answer silently, and the program carries on with it. Knowing the range of
the type in hand is part of choosing it.

\vfill\pagebreak
\section{Boolean type and Boolean algebra}
\label{sec:types:boolean}

A Boolean holds one of two values, \token{true} or \token{false}. That is the
whole type, and it is enough to build every decision a program makes: which
branch to take, whether to loop again, whether to stop. Ymir calls it
\token{bool} and stores it in one byte, in which \token{0b00000000} is
\token{false} and every other pattern is \token{true}.

\subsection{Primary operators}

Boolean algebra is the arithmetic of truth values. Where numerical algebra has
four basic operations, this one has three -- AND, OR and NOT -- and every other
logical operation can be written out of those three alone.

AND, or conjunction, is $\land$ in algebra and \token{&&} in Ymir. It takes two
operands and yields \token{true} only when both are \token{true}. OR, or
disjunction, is $\lor$ and \token{||}: it yields \token{true} when at least one
operand is. NOT, or negation, is $\lnot$ and \token{!}, takes one operand, and
returns the other value. In the tables that follow, \bot{} stands for
\token{false} and \top{} for \token{true}.

\begin{table}[h]
  \centering
  \adjustbox{scale=0.9, max width=\linewidth}{
  \begin{minipage}[t]{0.4\textwidth}
    \centering
    \begin{tabular}{p{.15\textwidth}|p{.15\textwidth}|p{.2\textwidth}|p{0.2\textwidth}}
      \toprule[0.6pt]
      $x$ & $y$ & $x \land y$ & $x \lor y$ \\
      \toprule[0.6pt]        
      \top & \top & \top & \top \\
      \top & \bot & \bot & \top \\
      \bot & \top & \bot & \top \\
      \bot & \bot & \bot & \bot \\
      \midrule[0.2pt]
    \end{tabular}
    \captionof{table}{AND, OR operators}
  \end{minipage}%
  \hspace{50pt}%
  \begin{minipage}[t]{0.3\textwidth}
    \centering
    \begin{tabular}{p{.15\textwidth}|p{.30\textwidth}}
      \toprule[0.6pt]
      $x$ & $\lnot x$ \\
      \toprule[0.6pt]        
      \top & \bot \\
      \bot & \top \\
      \midrule[0.2pt]
    \end{tabular}
    \captionof{table}{NOT operator}
  \end{minipage}
  }
%  \caption{\label{tab:types:boolean_operators} Boolean algebra operators}
\end{table}


\vspace{-10pt}%
\subsection{Secondary operators}

Three further operators are common enough to have names of their own --
implication, bidirectional implication and exclusive disjunction -- but none of
them is new: each rewrites into the three basic ones.

\begin{itemize}
\vspace{-5pt}%
\item Implication: $x \rightarrow y = \lnot x \lor y$.
\vspace{-5pt}%
\item Bidirectional implication : $x \leftrightarrow y = (\lnot x \lor y) \land (x \lor \lnot y)$ 
\vspace{-5pt}
\item Exclusive disjunction (or XOR) : $x \oplus y = (x \land \lnot y) \lor (\lnot x \land y)$
\end{itemize}
\vspace{-5pt}%

\begin{table}[h]
  \centering
  \adjustbox{scale=0.9, max width=\linewidth}{
    \begin{tabular}{p{.15\textwidth}|p{.15\textwidth}|p{.2\textwidth}|p{0.2\textwidth}|p{0.2\textwidth}}
      \toprule[0.6pt]
      $x$ & $y$ & $x \rightarrow y$ & $x \leftrightarrow y$ & $x \oplus y$ \\
      \toprule[0.6pt]        
      \top & \top & \top & \top & \bot \\
      \top & \bot & \bot & \bot & \top \\
      \bot & \top & \top & \bot & \top \\
      \bot & \bot & \top & \top & \bot \\
      \midrule[0.2pt]
    \end{tabular}
  }
  \caption{\label{tab:types:boolean_operators} Secondary operators}
\end{table}




The truth table above is what justifies those rewritings. Take implication. When
$x$ is true, the implication holds only if $y$ is true as well -- that is the
case everyday language recognises. When $x$ is false, the implication says
nothing about $y$ and holds either way, because implication runs one way only:
the truth of $y$ tells us nothing about $x$. Read as English, $\lnot x \lor y$
sounds like a different claim; row for row, it is the same one.

Ymir provides none of the three. The primary operators already express them, and
a language gains nothing by spelling them out.

\subsection{Laws on Boolean operators}

Boolean operations obey much the same laws as arithmetic ones, with $\lor$ in the
role of addition and $\land$ in the role of multiplication. Commutativity,
associativity and distributivity all carry over.

\textit{a)} Commutativity: the order of the operands does not affect the result of the
operation. This property applies to both AND and OR, for example $x \lor y =
y \lor x$.
\smallskip

\textit{b)} Associativity: the way operands are grouped does not affect the result of the
operation. This property applies to both AND and OR; for example, $x \land
(y \land z) = (x \land y) \land z$. However, associativity does not hold when
different operators are mixed. For instance, $x \land (y \lor z) \neq (x \land
y) \lor z$. To illustrate, let $x = \bot, y = \bot$ and $z = \top$, then
$x \land (y \lor z) = \bot$ while $(x \land y) \lor z = \top$.
\smallskip

\textit{c)} Distributivity: One operation can be spread across another. For instance
$x \land (y \lor z) = (x \land y) \lor (x \land z)$.

These laws are also where Boolean logic goes wrong most often. The analogy with
arithmetic extends to precedence -- \token{&&} binds tighter than \token{||},
just as multiplication binds tighter than addition -- and it is worth
parenthesising a Boolean expression even when the parentheses change nothing.

The classic mistake is to lean on commutativity and forget precedence, which
quietly regroups the operands. In the next listing \token{r1} and \token{r2} look
like rearrangements of one another and hold different values.

\begin{lstlisting}[style=coloredverbatim, label=lst:types:boolean_error, caption=Example of logic error when dealing with Boolean values, escapechar=@]
let x = false, y = false, z = true;

let r1 = x && y || z;
let r2 = (z || x) && y;

@\hb{assert(r1 == r2);}@ // Sorry, wrong. 
\end{lstlisting}

\subsection{Specificities of Boolean operators in Ymir}

Ymir compiles a Boolean operator into a branch, not into a calculation, and the
consequence is visible from the outside: the second operand is evaluated only if
it can still change the answer. A \token{false} on the left of \token{&&}, or a
\token{true} on the left of \token{||}, settles the result on its own, and the
right-hand side is skipped.

This buys two things. The skipped operand may be expensive -- a function call
rather than a variable -- so not evaluating it saves the work. More importantly,
it makes the left operand a guard for the right: later chapters lean on that
heavily.

In the listing below \token{foo} is never called. \token{x} is \token{true}, and
$\top \lor y$ is $\top$ whatever $y$ turns out to be. Swap the \token{||} on
line~11 for \token{&&} and \token{foo} does run -- unless \token{x} is also
changed to \token{false}, since $\bot \land y$ is $\bot$ regardless.

\begin{lstlisting}[style=coloredverbatim, caption=Example of boolean operation optimised out]
fn foo()-> bool {
  println("In foo");
  true
}

fn main()
  throws AssertError
{
  let mut x = true;

  assert(x || foo()); // not calling foo
  println("Exit");
}
\end{lstlisting}
\begin{lstlisting}[style=bashVerb, escapechar=@+]
@+\textcolor{teal!80}{alice@dev:\mtilde}@+$ gyc main.yr -o main
@+\textcolor{teal!80}{alice@dev:\mtilde}@+$ ./main
Exit
\end{lstlisting}

\token{assert} stops the program when its condition is \token{false}. A
function that calls it has to say so, with \token{throws AssertError} between
its parameters and its body, as on line~7. Until
Chapter~\ref{chap:error} explains why, treat that line as part of the formula.

So the commutativity that Boolean algebra guarantees does not survive contact
with Ymir. The \emph{value} of \token{a && b} and \token{b && a} is the same; what
each of them evaluates on the way there is not.

\vfill%
\pagebreak
\section{Floating point types}
\label{sec:types:floats}

Numbers with a decimal point are floating-point numbers. Where an integer type
trades away range -- it counts exactly, up to a limit -- a floating-point type
trades away exactness: it reaches enormous magnitudes, and pays for them with
values it can only approximate. How closely it approximates depends on how many
bits it is given.

There are four such types, named on the same pattern as the integers:
\token{f32}, \token{f64}, \token{f80}, and \token{fsize}, whose width is the
widest the machine supports.

\subsection{Floating point literals}

A floating-point literal is two runs of digits with a \token{.} between them and
no spaces anywhere. Underscores may separate digits, as in an integer literal.
Both parts are required: \token{10.0} is a literal, \token{10.} is not.

A literal with no suffix is an \token{f64}. The suffixes are \token{f} for
\token{f32}, \token{d} for \token{f64}, \token{l} for \token{f80} and \token{r}
for \token{fsize} -- the initials of \textit{float}, \textit{double},
\textit{long double} and \textit{real}, the names these widths have carried
since long before Ymir.

\begin{lstlisting}[style=coloredverbatim, label=lst:types:float_literal, caption=Example of floating point literals]
let x = 3.14_15_92f; // f32

let y = 2.71828d; // f64
let z = 2.71828; // f64

let w = 1.0l; // f80

let v = 1_000.0r; // fsize
let u = 1.888_999r; // fsize

println(x, " ", y, " ", z, " ", w, " ", v, " ", u);
\end{lstlisting}

\notebox{
  More advanced forms of floating-point literals are also available,
  such as hexadecimal notation and scientific notation. These representations
  are explained in detail in the advanced chapter on floating-point types (see
  Section~\ref{sec:scalars:floats}).
}

\subsection{Memory representations}

Floating-point numbers follow the IEEE-754 standard, which splits every value
into three fields: a sign, an exponent, and a mantissa -- also called the
significand. How many bits each field gets depends on the type.

\begin{table}[H]
  \centering
  \begin{tabular}{p{.15\textwidth}|p{.25\textwidth}|p{0.20\textwidth}|p{0.20\textwidth}}
    \toprule[0.6pt]
    \textbf{Type} & \textbf{Sign (bits)} & \textbf{Exponent (bits)} & \textbf{Mantissa (bits)} \\
    \toprule[0.6pt]
    \token{f32}   & 1 & 8  & 23 \\
    \token{f64}   & 1 & 11 & 52 \\
    \token{f80}   & 1 & 15 & 64\\
    \token{fsize} & 1 & arch. dependent & arch. dependent \\
    \bottomrule[0.2pt]
  \end{tabular}
  \caption{Bit layout of floating-point types in Ymir (IEEE-754 representation).}
\end{table}


The value of a floating-point number is computed using the following formula,
where the sign is either $+1$ or $-1$, and the bias is defined as $bias =
2^{e-1} - 1$ with $e$ representing the number of bits used for the exponent. For
instance, in the \token{f32} type, $e = 8$, which gives a bias of $127$.

\begin{equation*}
  value = sign \times mantissa \times 2^{(exponent - bias)}
\end{equation*}

The bias is what lets the exponent go negative without a sign bit of its own: an
exponent of $0$ is stored as $127$, $1$ as $128$, $-1$ as $126$. Positive and
negative exponents end up spread evenly on either side of the stored midpoint,
which keeps the hardware simple. For \token{f32} the exponent runs from about
$-126$ to $127$.

\subsubsection{Example}

Take $87.25$ as an \token{f32}, and work it through field by field.

First, binary. The integer part $87$ is $1010111_2$. The fractional part $0.25$
is $0.01_2$, since $0.25 = 0 \times 2^{-1} + 1 \times 2^{-2} + 0 \times 2^{-3} +
\ldots$ Together they give $1010111.01_2$.

Next, normalisation into the form $1.mantissa \times 2^{exponent}$: shifting the
point six places left gives $1.01011101 \times 2^{6}$. The exponent is therefore
$6$, and since the number is positive the sign bit is $0$.

Then the bias. For \token{f32} it is $127$, so the exponent field holds $6 + 127
= 133$, which on eight bits is $10000101_2$. The mantissa field holds the digits
after the point of $1.01011101$ -- the leading $1$ is implied and not stored --
padded with zeros out to the 23 bits the format allots it.
Table~\ref{tab:types:float_mem_repr_ex} assembles the three.

\begin{table}[h!]
  \centering
  %% The three fields are bit patterns, not quantities: they are set in the
  %% monospaced face like every other bit pattern in the book. As maths they
  %% came out in wide-spaced italics, and the 23-bit mantissa then overran its
  %% column.
  \begin{tabular}{p{.10\textwidth}|p{.22\textwidth}|p{0.42\textwidth}}
    \toprule[0.6pt]
    \textbf{Sign} & \textbf{Exponent} & \textbf{Mantissa} \\
    \toprule[0.6pt]
    \token{0} & \token{10000101} & \token{01011101000000000000000} \\
    \bottomrule[0.2pt]
  \end{tabular}
  \caption{\label{tab:types:float_mem_repr_ex} Memory representation of $87.25$ in a \token{f32}}
\end{table}

\subsection{Stepping}

Between any two decimal numbers there are infinitely many others, and a fixed
number of bits cannot name them all. A floating-point type therefore represents
some values exactly and, for everything in between, silently substitutes the
nearest value it does have. The gaps between the values it has are called
\textit{stepping}.

Listing~\ref{lst:types:stepping_example} shows what that costs. It subtracts
\token{x} from \token{z} to get \token{y}, then adds \token{x} back to \token{y}
to get \token{w}. Arithmetic says \token{w} equals \token{z}. The program says
otherwise: each operation landed on the nearest representable value, and the two
roundings did not cancel. The example uses \token{f32}, and \token{f64} happens
to survive it -- but only this one. Wider types have finer steps, not no steps.

\begin{lstlisting}[style=coloredverbatim, label=lst:types:stepping_example, caption=Example of floating point stepping]
let mut z = 1.989328f; // 'f' suffix to create a f32 literal
let mut x = 0.98769778678f;
let y = z - x;

let w = x + y;

println("(X : ", x, ", Y : ", y, ")");
println("(Z : ", z, ", W : ", w, ")");
println("Z == W : ", z == w);
\end{lstlisting}
\vspace{-10pt}%
\begin{lstlisting}[style=bashVerb]
(X : 0.987698, Y : 1.00163)
(Z : 1.98933, W : 1.98933)
Z == W : false
\end{lstlisting}

\cautionbox{the variable \token{z} and \token{x} are mutable in the
  Listing~\ref{lst:types:stepping_example} to prevent compile time
  optimization and ensure the arithmetic operation are indeed executed at
  runtime.}

\subsection{Special values}

Some bit patterns are reserved. IEEE-754 sets aside particular combinations of
exponent and mantissa to mean infinity, zero, and NaN -- Not a Number, the result
of an operation that has no answer at all.

\begin{table}[H]
  \centering
  \begin{tabular}{p{.15\textwidth}|p{.15\textwidth}|p{0.10\textwidth}|p{0.10\textwidth}|p{0.30\textwidth}}
    \toprule[0.6pt]
    \textbf{Special Value} & \textbf{Exponent} & \textbf{Mantissa} & \textbf{Sign} & \textbf{Meaning} \\
    \toprule[0.6pt]
    $+0$ / $-0$ & $0$ & $0$ & $0$ / $1$ & Zero with sign \\
    $+\infty$ / $-\infty$ & all bits set to 1 & $0$ & $0$ / $1$ & Overflow or division by zero \\
    NaN & all bits set to 1 & $\neq 0$ & $0$ / $1$ & Invalid of undefined result \\
    Subnormal & $0$ & $\neq 0$ & $0$ / $1$ & Very small numbers close to zero \\
    \bottomrule[0.2pt]
  \end{tabular}
  \caption{Special value representable using IEEE-754 standard}
\end{table}


\vfill\pagebreak
\section{Character and String types}
\label{sec:types:chars_strings}

A character type holds one textual symbol -- a letter, a digit, a punctuation
mark. Put characters in sequence and the result is a string, which is how a
program holds anything longer than a symbol: a word, a line, a file.

\subsection{Character literal}
\label{sec:types:char_literal}

A character literal is a symbol between single quotes, \token{'}. Three types
hold one: \token{c8}, \token{c16} and \token{c32}, for UTF-8, UTF-16 and UTF-32
respectively. The encoding is what decides how many bits a symbol takes, and once
strings enter the picture that choice stops being a detail --
Section~\ref{sec:types:encoding} is about little else.

\begin{lstlisting}[style=coloredverbatim, label=lst:types:char_literal, caption=Example of character literals, escapechar=@]
let x = 'A';
let y = '@\ensuremath{\color{black!30!applegreen} \texttt{\Omega{}}}@'c16;
let z = '@\emoji{slightly-smiling-face}@'c32;

println(x, " ", y, " ", z);
\end{lstlisting}


Listing~\ref{lst:types:char_literal} stores three of them. A suffix names the
type, as it does for numbers: \token{'Ω'c16} is a \token{c16} and
\token{'🙂'c32} a \token{c32}. An unsuffixed literal such as \token{'A'} is
always a \token{c8}, with no exception made for wide characters.

Which is why \token{'🙂'} on its own does not compile. That symbol needs a full
32 bits, and an unsuffixed literal offers it 8.

\begin{lstlisting}[style=coloredverbatimError, label=lst:types:char_encode_error, escapechar=@]
let x = '@\hb{\emoji{slightly-smiling-face}}@'; // Error : malformed literal, number of c8 is 4 
\end{lstlisting}%

Some characters cannot be typed inside a literal: a line break would end the
line, and a quote would end the literal. These are written as escape sequences
instead, always introduced by a backslash. Ymir accepts the ones below, and all
of them except \tokennolst{\textbackslash{}u\{code\}} -- whose width depends on
the code point it names -- fit in a single UTF-8 character.

\begin{center}
  \begin{tabular}{ll}
    Value & Content\\[0pt]
    \hline
    \tokennolst{\textbackslash{}a} & Alert beep, (Bell)\\[0pt]
    \tokennolst{\textbackslash{}b} & Backspace\\[0pt]
    \tokennolst{\textbackslash{}f} & Page break\\[0pt]
    \tokennolst{\textbackslash{}n} & New line\\[0pt]
    \tokennolst{\textbackslash{}r} & Carriage return\\[0pt]
    \tokennolst{\textbackslash{}t} & Horizontal tab\\[0pt]
    \tokennolst{\textbackslash{}v} & Vertical tab\\[0pt]
    \tokennolst{\textbackslash{}\textbackslash{}} & Backslash\\[0pt]
    \tokennolst{\textbackslash{}'} & Apostrophe\\[0pt]
    \tokennolst{\textbackslash{}"} & Double quotation mark\\[0pt]
    \tokennolst{\textbackslash{}u\{code\}} & Unicode\\[0pt]
  \end{tabular}
  \captionof{table}{\label{tab:types:escape_chars} Escape characters}
\end{center}

The last of them, \tokennolst{\textbackslash{}u\{code\}}, is the way to write a
character by its Unicode code point rather than by the symbol itself -- useful
when the symbol is hard to type, or hard to tell apart from another. The code
point of \token{'🙂'} is \token{0x1F642}, so it may equally be written
\token{'\u\{0x1F642\}'c32}.

\subsection{String type}

A string holds text: characters laid out in memory one after another, in some
encoding. A single character type is fixed in width, but a string is not -- what
it costs in memory depends on which symbols are in it.

\subsection{String literal}

A string literal is a run of symbols between double quotes, \token{"}. Its type
is written in brackets, and reads as an array of characters: \token{\[c8\]} is a
UTF-8 string of any length, \token{\[c16 ; 5\]} a UTF-16 string of exactly five.
The notation states the encoding and the length constraint in one place. Arrays
and slices are the subject of a later chapter; all that matters here is that both
store their elements consecutively in memory.

String literals take suffixes too -- \token{s8}, \token{s16} and \token{s32} --
and default to UTF-8 without one.

\begin{lstlisting}[style=coloredverbatim, label=lst:types:string_literal, caption=Example of string literals, escapechar=@]
let greetings = "Hello World!";
let utf8_unicode = "@\korean{{\color{black!30!applegreen} 안녕하세요 세계}}@";
let utf32_unicode = "@\japan{{\color{black!30!applegreen} こんにちは 世界}}@"s32;

println(greetings, " ", utf8_unicode, " ", utf32_unicode);
\end{lstlisting}

\subsection{A small word on encoding}
\label{sec:types:encoding}

A character may take anything from one to four bytes, as
Section~\ref{sec:types:char_literal} noted. The characters of the ASCII table
take one, which is why they sit in a \token{c8}, or in a \token{[c8]} string,
with nothing clever going on. The \token{greetings} variable of
Listing~\ref{lst:types:string_literal} is one of these: plain ASCII, one byte per
character, laid out in order, starting with \token{'H'} -- which is $72$.

\begin{figure}[h]
  \centering
  \begin{tikzpicture}
    \draw (0, 0) -- coordinate (A) (0, 0);
    \draw (-6, -1) -- coordinate (B) (-6, -1);
    
    \filldraw ($(A)$) node[align=center] {\color{black!30!applegreen}"Hello World!"};

    \draw[-] ($(B)+(0, 0)$) rectangle ($(B)+(12, -0.5)$);
    \draw[-] ($(B)+(1, 0)$) -- ($(B)+(1, -0.5)$);
    \draw[-] ($(B)+(2, 0)$) -- ($(B)+(2, -0.5)$);
    \draw[-] ($(B)+(3, 0)$) -- ($(B)+(3, -0.5)$);
    \draw[-] ($(B)+(4, 0)$) -- ($(B)+(4, -0.5)$);
    \draw[-] ($(B)+(5, 0)$) -- ($(B)+(5, -0.5)$);
    \draw[-] ($(B)+(6, 0)$) -- ($(B)+(6, -0.5)$);
    \draw[-] ($(B)+(7, 0)$) -- ($(B)+(7, -0.5)$);
    \draw[-] ($(B)+(8, 0)$) -- ($(B)+(8, -0.5)$);
    \draw[-] ($(B)+(9, 0)$) -- ($(B)+(9, -0.5)$);
    \draw[-] ($(B)+(10, 0)$) -- ($(B)+(10, -0.5)$);
    \draw[-] ($(B)+(11, 0)$) -- ($(B)+(11, -0.5)$);
    \draw[-] ($(B)+(12, 0)$) -- ($(B)+(12, -0.5)$);

    \filldraw ($(B)+(0.5, -0.25)$) node[align=center] {$72$};
    \filldraw ($(B)+(1.5, -0.25)$) node[align=center] {$101$};
    \filldraw ($(B)+(2.5, -0.25)$) node[align=center] {$108$};
    \filldraw ($(B)+(3.5, -0.25)$) node[align=center] {$108$};
    \filldraw ($(B)+(4.5, -0.25)$) node[align=center] {$111$};
    \filldraw ($(B)+(5.5, -0.25)$) node[align=center] {$32$};
    \filldraw ($(B)+(6.5, -0.25)$) node[align=center] {$87$};
    \filldraw ($(B)+(7.5, -0.25)$) node[align=center] {$111$};
    \filldraw ($(B)+(8.5, -0.25)$) node[align=center] {$114$};
    \filldraw ($(B)+(9.5, -0.25)$) node[align=center] {$108$};
    \filldraw ($(B)+(10.5, -0.25)$) node[align=center] {$100$};    
    \filldraw ($(B)+(11.5, -0.25)$) node[align=center] {$33$};

    \draw[->] ($(A)+(-1, -0.2)$) -- ($(B)+(0.5, 0.1)$);
    \draw[->] ($(A)+(-0.75, -0.2)$) -- ($(B)+(1.5, 0.1)$);
    \draw[->] ($(A)+(-0.6, -0.2)$) -- ($(B)+(2.5, 0.1)$);
    \draw[->] ($(A)+(-0.45, -0.2)$) -- ($(B)+(3.5, 0.1)$);
    \draw[->] ($(A)+(-0.3, -0.2)$) -- ($(B)+(4.5, 0.1)$);
    \draw[->] ($(A)+(-0.1, -0.2)$) -- ($(B)+(5.5, 0.1)$);
    \draw[->] ($(A)+(0.1, -0.2)$) -- ($(B)+(6.5, 0.1)$);

    \draw[->] ($(A)+(0.4, -0.2)$) -- ($(B)+(7.5, 0.1)$);
    \draw[->] ($(A)+(0.55, -0.2)$) -- ($(B)+(8.5, 0.1)$);
    \draw[->] ($(A)+(0.75, -0.2)$) -- ($(B)+(9.5, 0.1)$);
    \draw[->] ($(A)+(0.9, -0.2)$) -- ($(B)+(10.5, 0.1)$);
    \draw[->] ($(A)+(1.1, -0.2)$) -- ($(B)+(11.5, 0.1)$);    
    
  \end{tikzpicture}

  \caption{\label{fig:types:hello_string_figure} Encoding of a simple UTF-8 string value in memory}
\end{figure}



Everything outside ASCII goes through Unicode, which assigns a code point to
every character in every writing system. The three encodings differ in how they
store those code points, and the trade-off is always the same: simplicity against
space.

UTF-32 takes the simple end. Every character gets four bytes, whether it needs
them or not, and the surplus is zeros. Nothing is ever ambiguous, and a good deal
of memory is wasted.



The following Japanese text \colorbox{background!90!gray!60}{\japan{\color{black!30!applegreen} "こんにちは"}}
is composed of five characters, each represented as a UTF‑32
Unicode code point: \token{0x3053 0x3093 0x306b 0x3061 0x306f}. The first
character, \colorbox{background!90!gray!60}{\japan{\color{black!30!applegreen} こ}}, has the Unicode code point
\token{0x3053}, which requires only 14 bits. This means it could be stored
within a \token{c16} type. However, under UTF‑32 encoding, every character is
allocated a full 4 bytes, leaving two bytes filled with zeros. The same
applies to each of the five characters in the word, resulting in a total of 20
bytes used, even though only 10 bytes contain meaningful information. On the
other hand, the character \token{🙂} has the Unicode code point
\token{0x1f642}, which requires 17 bits and therefore no longer fits inside a
\token{c16}. As a result, UTF‑16 and UTF‑32 encodings of that character take
the same number of bytes to store it in memory.



UTF‑8 encodes Unicode characters by using a variable number of bytes for each
symbol. The first byte contains information that indicates how many subsequent
bytes belong to the same character. For example, in the mixed English and Korean
text \colorbox{background!90!gray!60}{\color{black!30!applegreen}{"\texttt{Hello is }\korean{안녕하세요}"}}, the segment \colorbox{background!90!gray!60}{\color{black!30!applegreen}{"\texttt{Hello is }"}} uses one
byte per character, while the segment \colorbox{background!90!gray!60}{\korean{\color{black!30!applegreen}{"안녕하세요"}}} requires three bytes
for each character. The UTF‑8 encoding of the entire string therefore occupies
24 bytes of memory, compared to 56 bytes in UTF‑32 and 28 bytes in UTF‑16.


The saving has a catch, and it is the important part of this section: the length
of a UTF-8 string is a count of bytes, not of characters.
\colorbox{background!90!gray!60}{\korean{\color{black!30!applegreen}{"안녕하세요"}}} is
five characters and fifteen bytes, and it is fifteen that a \token{\[c8\]}
reports. UTF-16 has the same problem in milder form, its characters taking two
bytes or four depending on the code point.

\begin{lstlisting}[style=coloredverbatim, label=lst:types:encoding_example, caption=Encoding consequence on array length, escapechar=@]
let utf8_str = "@\korean{\color{black!30!applegreen} 안녕하세요 세계}@";
let utf16_str = "@\korean{\color{black!30!applegreen} 안녕하세요 세계}@"s16;
let utf32_str = "@\korean{\color{black!30!applegreen} 안녕하세요 세계}@"s32;

println(utf8_str.len); // 22, or 22 bytes
println(utf16_str.len); // 8, or 8 * 2 => 16 bytes
println(utf32_str.len); // 8, or 8 * 4 => 32 bytes
\end{lstlisting}

\notebox{ For all these encoding reasons, it becomes clear why the character
  \emoji{slightly-smiling-face} cannot be represented using a single \token{c8}. In fact, the
  length of \token{"🙂"} is reported as 4, since UTF‑8 requires four bytes to
  encode this symbol. Encoding is not a trivial topic. It will be explored in
  greater detail in a later chapter. What is important to understand here is
  that the apparent length of a string does not always correspond to the number
  of characters it contains. When in doubt, converting the string into UTF‑32
  will reveal the true character count, as UTF‑32 consistently encodes each
  character with a fixed length of four bytes. }

\begin{lstlisting}[style=coloredverbatim, label=lst:types:converting_encoding, caption=String encoding conversion, escapechar=@]
use std::conv; // to import 'to'

let utf8_str = "@\korean{\color{black!30!applegreen} 안녕하세요 세계}@";

let utf32_str = to!{[c32]}(utf8_str); // converting UTF8 into UTF32

println(utf8_str.len); // 22, bytes
println(utf32_str.len); // 8, characters
\end{lstlisting}

\vfill%
\pagebreak

\section*{Exercises}
\markright{\MakeUppercase{Exercises}}%
\subsection*{Exercise 1}

Write a Ymir program that declares a mutable \token{i32} variable named
\token{score}, initialized to \token{5}, and prints it. Then, update
\token{score} so that it holds its own square, and print it again.

\subsection*{Exercise 2}

What does the following program print? Work out the result by hand using the
precedence rules from Table~\ref{tab:types:operator_precedence}, then check
your answer by compiling and running the program.

\begin{lstlisting}[style=coloredverbatim]
use std::io;

fn main() {
  let a = 2;
  let b = 5;
  let c = 3;

  let r = a + b * c - a;
  println("R : ", r);
}
\end{lstlisting}

\subsection*{Exercise 3}

The following program contains two errors, one for each kind of mistake this
chapter introduced.

\begin{lstlisting}[style=coloredverbatimError, escapechar=@, caption=Ymir program with errors]
use std::io;

fn main() {
  let x: i32 = 10;
  @\hb{x = 20;}@

  let y: f32 = @\hb{"text"}@;

  println("X : ", x);
  println("Y : ", y);
}
\end{lstlisting}

Correct both errors and compile the program with \token{gyc} to test your
corrections.

\subsection*{Exercise 4}

What does the following program print? \token{u8} is an unsigned 8-bit
integer type.

\begin{lstlisting}[style=coloredverbatim]
use std::io;

fn main() {
  let mut x: u8 = 250;

  x += 10;
  println("X : ", x);
}
\end{lstlisting}

\subsection*{Exercise 5}

What does the following program print? For each of the five variables, state
which floating-point or integer type it holds, and explain why.

\begin{lstlisting}[style=coloredverbatim]
use std::io;

fn main() {
  let a = 12;
  let b = 12u64;
  let c = 3.14;
  let d = 3.14f;
  let e = 1_000_000;

  println("A : ", a);
  println("B : ", b);
  println("C : ", c);
  println("D : ", d);
  println("E : ", e);
}
\end{lstlisting}

\subsection*{Exercise 6}

What does the following program print? The word \textit{café} has four
characters. Explain why the printed value is not 4.

\begin{lstlisting}[style=coloredverbatim]
use std::io;

fn main() {
  let word = "café";

  println(word.len);
}
\end{lstlisting}

\subsection*{Exercise 7}

What does the following program print? List, in order, which calls to
\token{check} are actually made, and explain why the other ones are skipped.

\begin{lstlisting}[style=coloredverbatim]
use std::io;

fn check(name: [c8], value: bool)-> bool {
  println("Checking ", name);
  value
}

fn main() {
  let mut ok = true;

  if check("a", false) && check("b", true) {
    ok = false;
  }

  println("OK : ", ok);
}
\end{lstlisting}

\vfill%
\pagebreak

\forceevenpage{}

\section*{Solutions}
\markright{\MakeUppercase{Solutions}}%
\subsection*{Exercise 1}

\begin{lstlisting}[style=coloredverbatim, caption=Solution for exercise 1]
use std::io;

fn main() {
  let mut score = 5;
  println("Score : ", score);

  score = score * score;
  println("Score : ", score);
}
\end{lstlisting}

\subsection*{Exercise 2}

The multiplication \token{b * c} is evaluated first, since \token{*} has
higher precedence than \token{+} and \token{-}. The remaining two operators
share the same precedence, so they are evaluated left to right:
\token{(a + (b * c)) - a}, that is \token{(2 + 15) - 2 = 15}.

\begin{lstlisting}[style=bashVerb]
R : 15
\end{lstlisting}

\subsection*{Exercise 3}

The variable \token{x} is declared with \token{let} but not \token{let mut},
so it is a constant and cannot be reassigned. The variable \token{y} is
declared with the type \token{f32}, but the value on the right-hand side of
the \token{=} sign is a string literal, not a floating-point literal — the
types of the \textit{lvalue} and the \textit{rvalue} do not match.

\begin{lstlisting}[style=coloredverbatim, escapechar=@, caption=Solution for exercise 3]
use std::io;

fn main() {
  let @\hcb{mut}@ x: i32 = 10;
  x = 20;

  let y: f32 = @\hcb{2.5}@;

  println("X : ", x);
  println("Y : ", y);
}
\end{lstlisting}

\subsection*{Exercise 4}

Adding \token{10} to \token{250} exceeds the maximum value representable by a
\token{u8} (\token{255}), so the result overflows and wraps around:
\token{250 + 10 = 260}, and \token{260 mod 256 = 4}.

\begin{lstlisting}[style=bashVerb]
X : 4
\end{lstlisting}

\subsection*{Exercise 5}

\begin{lstlisting}[style=bashVerb]
A : 12
B : 12
C : 3.14
D : 3.14
E : 1000000
\end{lstlisting}

\token{a} carries no suffix, so it defaults to \token{i32}. \token{b} carries
the explicit suffix \token{u64}, so it is a 64-bit unsigned integer.
\token{c} carries no suffix, so it defaults to \token{f64}. \token{d} carries
the suffix \token{f}, so it is an \token{f32}. \token{e} is again an
\token{i32}; the underscores only group the digits for readability and do not
change the value or the type.

\subsection*{Exercise 6}

\begin{lstlisting}[style=bashVerb]
5
\end{lstlisting}

The word has four characters, but \token{é} is not part of the ASCII table
and requires two bytes to be encoded in UTF‑8, while every other character of
\textit{café} requires only one. The default string type is \token{[c8]},
whose \token{.len} reports a number of bytes, not a number of characters, so
the printed length is $3 + 2 = 5$.

\subsection*{Exercise 7}

\begin{lstlisting}[style=bashVerb]
Checking a
OK : true
\end{lstlisting}

Only \token{check("a", false)} is called. The \token{\&\&} operator is the
Boolean AND: as soon as its left operand is \token{false}, the whole
expression is necessarily \token{false}, so the right operand,
\token{check("b", true)}, is never evaluated, and \token{"Checking b"} is
never printed. Since the condition of the \token{if} is \token{false}, its
body is skipped and \token{ok} keeps its initial value \token{true}.

